\documentclass[11pt]{article}
\usepackage[a4paper,margin=27mm]{geometry}
\usepackage{amsmath,amssymb,amsthm,mathtools}
\usepackage[T1]{fontenc}
\usepackage[hidelinks]{hyperref}
\newtheorem{theorem}{Theorem}[section]
\newtheorem{lemma}[theorem]{Lemma}
\newtheorem{proposition}[theorem]{Proposition}
\newtheorem{corollary}[theorem]{Corollary}
\theoremstyle{definition}
\newtheorem{definition}[theorem]{Definition}
\theoremstyle{remark}
\newtheorem{remark}[theorem]{Remark}
\newcommand{\NN}{\mathbb N}
\newcommand{\ZZ}{\mathbb Z}
\newcommand{\RR}{\mathbb R}
\newcommand{\CC}{\mathbb C}
\newcommand{\cS}{\mathcal S}
\newcommand{\htp}{\operatorname{ht}}
\newcommand{\one}{\mathbf 1}
\setlength{\emergencystretch}{3em}
\hypersetup{pdftitle={Infinite paths in the composite-restricted visible lattice},pdfauthor={Alex Chengyu Li}}
\title{Infinite paths in the composite-restricted visible lattice}
\author{Alex Chengyu Li\\\small Independent Researcher}
\date{6 September 2026}

\begin{document}
\maketitle
\begin{abstract}
We consider the nearest-neighbour graph on pairs of integers greater than
one that are coprime and have at least one composite coordinate. We prove
that this graph contains an infinite simple path, answering Erd\H{o}s
Problem~1212 affirmatively. More precisely, for almost every slope in a
fixed interval away from the diagonal, there is a ray with that limiting
coordinate ratio. The proof first obtains many pairwise coprime composite
rows free of small prime factors in a short band. A polynomial Jacobsthal bound supplies
composite supporting columns, so a bad crossing through the band has a
subcrossing of polylogarithmic width. Its large common prime divisors then
force an integer interpolation polynomial of small height to vanish at
one of its vertices. A summable estimate for the exceptional directions of
all such polynomials permits one direction to be fixed at every sufficiently
large scale. Planar crossing duality and explicit overlapping rectangles
produce an unbounded connected subgraph along this direction, from which
an infinite simple path is extracted.
\end{abstract}

\section{The graph and the main theorem}

Let
\begin{equation}\label{eq:safe-set}
 \cS=\{(x,y)\in\NN^2:x,y>1,\ \gcd(x,y)=1,
                 \text{ at least one of }x,y\text{ is composite}\}.
\end{equation}
Two vertices are adjacent when one coordinate changes by one and the other
does not change. A ray is a sequence of distinct vertices with consecutive
vertices adjacent. Every ray in this locally finite graph leaves every
finite set.

Erd\H{o}s asked whether this graph has a path going to infinity
\cite[p.~114]{Erdos1980}. The composite-coordinate condition is important:
Stewart's construction for the unrestricted visible lattice uses vertices
whose two coordinates are prime. Vardi studied the infinite component of
the unrestricted visible lattice and its asymptotic density \cite{Vardi1999}.
Random-coset versions of coprime percolation have also been studied
\cite{LeFournLiuMartineau2025}. Here every path is constructed in the
deterministic graph~\eqref{eq:safe-set}.

\begin{theorem}\label{thm:directions}
For Lebesgue-almost every $\alpha\in(4/3,5/3)$ there is a ray
$((x_t,y_t))_{t\geq0}$ in $\cS$ such that
\[
                         \frac{x_t}{y_t}\longrightarrow\alpha.
\]
In particular, the graph on $\cS$ has an infinite path.
\end{theorem}

For example, $(4,9),(5,9),(5,8)$ is a finite path in $\cS$.
The vertex $(4,7)$, however, is isolated: its horizontal neighbours have two
prime coordinates, and its vertical neighbours are nonvisible. Thus the
theorem asserts an infinite component in a graph that also contains
isolated vertices.

The central step is to associate a small integer polynomial to a long
\emph{bad} crossing, where a point with both coordinates greater than one
is bad when it is outside $\cS$. An integer is called $z$-rough when all
its prime factors exceed $z$. We obtain a short band containing many
rough composite rows. Safe vertical supports localize a bad subcrossing
in that band. At every selected row, a bad vertex has a large common prime
divisor; the divisors are distinct. An interpolation polynomial in the
short offsets then has an evaluation divisible by their product. This
product exceeds the size of a nonzero evaluation, so the evaluation is zero.
The polynomial may depend on the crossing. We handle this dependence by
excluding the directions associated with \emph{all} polynomials of the
required degree and height at once. Finally, overlapping rectangles in a
single narrowing corridor connect the crossings at successive scales.

\section{Notation and preliminary facts}

An integer interval $[a,b]_{\ZZ}$ contains both endpoints and has
$b-a+1$ sites. An integer rectangle is a product of two such intervals.
A left-right crossing is a path inside the rectangle joining its left and
right sides; top-bottom crossings are defined similarly. Safe paths use
the four nearest neighbours. A bad star path may also use a diagonal step,
so the maximum of the two coordinate changes is one.
We write $\omega(n)$ for the number of distinct prime divisors of $n$,
$\mu(n)$ for the Moebius function, and $p_k$ for the $k$th prime,
with $p_1=2$.

\begin{lemma}[Finite crossing facts]\label{lem:planar}
In an integer rectangle, a safe left-right crossing exists if and only if
there is no bad top-bottom star crossing. Furthermore, any nearest-neighbour
left-right and top-bottom crossings of the same rectangle have a common
vertex.
\end{lemma}
\begin{proof}
These are the finite four/eight-neighbour crossing alternatives. One way
to obtain the alternative is to join an auxiliary left side to every
safe site reachable from the left, and trace the boundary on the right of
this reachable set. If the set does not reach the right side, the boundary
has a segment joining the upper and lower sides. The lattice sites just
outside it are bad; successive such sites share an edge or a corner and
give the required star crossing. Repetitions can be removed. Conversely,
a left-right lattice path separates the upper and lower sides after its
ends are extended outside the rectangle. A top-bottom star crossing must
meet it. A diagonal edge between lattice sites cannot cross a horizontal
or vertical lattice edge in its interior without a common lattice endpoint,
so such an intersection would be a site of both colours. This is impossible.
The same separation argument for two nearest-neighbour crossings proves
the last assertion.
\end{proof}

We use the standard estimates
\begin{equation}\label{eq:prime-estimates}
 \prod_{p\leq z}(1-p^{-1})\sim\frac{e^{-\gamma}}{\log z},\qquad
 \sum_{z<p\leq2z}\frac1p\sim\frac{\log2}{\log z},\qquad
 \sum_{p\leq z}\frac1p=\log\log z+O(1),
\end{equation}
and $p_k=O(k\log(k+1))$. These follow from the prime number theorem,
Mertens' theorem and partial summation; see \cite{RosserSchoenfeld1962}.
We use Vaughan's uniform Jacobsthal bound \cite[p.~329]{Vaughan1977}:
for any set $Q$ of $k$ primes, every interval of
\begin{equation}\label{eq:jacobsthal}
                         O((k+1)^2\log^4(k+2))
\end{equation}
consecutive integers contains an integer coprime to $\prod_{q\in Q}q$.
The constant is absolute and the primes in $Q$ are arbitrary.

Throughout the proof $N$ is sufficiently large, and we write
\begin{equation}\label{eq:parameters}
 \ell=\log N,\quad L=\log\ell,\quad u=1000L,\quad z=N^{1/u},\quad
 d=\lceil4u\rceil,\quad r=\binom{d+2}{2}-1.
\end{equation}
The numerical parameter comparisons and counting estimates hold above a
single absolute lower bound for $N$, uniformly in interval positions,
chosen rows and bad paths. Later, the scale after which a fixed direction
avoids the exceptional sets is allowed to depend on that direction.

\section{Rough composite rows in short bands}

\begin{lemma}\label{lem:rough-composites}
Every integer interval $J\subset[N/2,4N]$ with $T=|J|\geq N^{9/10}$
contains at least
\begin{equation}\label{eq:rough-count}
                          \frac{T}{100(\log z)^2}
\end{equation}
composite integers whose prime factors all exceed $z$.
\end{lemma}
\begin{proof}
Put $P=\prod_{p\leq z}p$, $V=\prod_{p\leq z}(1-1/p)$, and
$\mathcal Q=\{q\text{ prime}:z<q\leq2z\}$, with
$S=\sum_{q\in\mathcal Q}q^{-1}$.
For an integer $h$ coprime to $P$, let
\[
                  A_h=\#\{t\in J:h\mid t,\ \gcd(t,P)=1\}.
\]
Choose an odd integer $m=100L+O(1)$. Bonferroni's inequalities at orders
$m$ and $m+1$ bound $A_h$ below and above by
\[
 \sum_{\substack{a\mid P\\\omega(a)\leq m}}\mu(a)
                       \#\{t\in J:ha\mid t\},\qquad
 \sum_{\substack{a\mid P\\\omega(a)\leq m+1}}\mu(a)
                       \#\{t\in J:ha\mid t\},
\]
respectively. Each count is $T/(ha)+O(1)$, with absolute error at most one.
The gap between the corresponding truncated main coefficients and $V$
is at most the elementary symmetric sum of order $m+1$ in the numbers
$1/p$, $p\leq z$. By~\eqref{eq:prime-estimates}, their sum is at most
$2L$ eventually. Thus the coefficient error is at most
\begin{equation}\label{eq:bonferroni-error}
 \varepsilon_N=\left(\frac{2eL}{m+1}\right)^{m+1}
                              \leq\ell^{-200}.
\end{equation}
The number of divisors retained at either order is at most
$E_N=(m+2)z^{m+1}$. We have obtained, uniformly in $h$ coprime to $P$,
\begin{equation}\label{eq:ah-count}
 \frac{T}{h}(V-\varepsilon_N)-E_N
 \leq A_h\leq
 \frac{T}{h}(V+\varepsilon_N)+E_N.
\end{equation}

Apply the first two inclusion-exclusion terms to divisibility by primes
in $\mathcal Q$. The number of integers counted by at least one of these
events is at least
\[
             \sum_{q\in\mathcal Q}A_q
                -\sum_{\substack{q,q'\in\mathcal Q\\q<q'}}A_{qq'}.
\]
Using~\eqref{eq:ah-count} and
$\sum_{q<q'}(qq')^{-1}\leq S^2/2$, this is at least
\begin{equation}\label{eq:rough-union}
 T V(S-S^2/2)-T\varepsilon_N(S+S^2/2)-O(N^{1/5}).
\end{equation}
Indeed, the summed integer-count remainder is at most
$4(m+2)z^{m+3}\leq N^{1/5}$ eventually, because
$(m+3)/u\longrightarrow1/10$.

By~\eqref{eq:prime-estimates}, the main coefficient in
\eqref{eq:rough-union} is asymptotic to
$e^{-\gamma}\log2/(\log z)^2$.
Both errors are $o(T/(\log z)^2)$ uniformly for $T\geq N^{9/10}$.
This proves~\eqref{eq:rough-count}. Finally, every counted integer exceeds
$N/2>2z$ and has a divisor $q\in\mathcal Q$, so is composite.
\end{proof}

\begin{lemma}\label{lem:band}
Set $B=\lceil10^4\ell^2\rceil$. Every interval $J$ as in
Lemma~\ref{lem:rough-composites} contains a block $J_0$ of exactly $B$
consecutive integers and distinct composite integers $c_1,\ldots,c_r\in J_0$
whose prime factors exceed $z$. The $c_i$ are pairwise coprime.
\end{lemma}
\begin{proof}
Partition $J$ into full blocks of $B$ integers and a final shorter block.
Removing the latter loses at most $B=o(T/(\log z)^2)$ selected integers.
By averaging, one full block contains at least
\[
             \frac{B}{200(\log z)^2}\geq50u^2
\]
of them. Since $r=8u^2+O(u)$, this is sufficient.
Also $B<z$ eventually. A common prime factor of two selected integers
would exceed $z$ but divide a nonzero difference smaller than $B$, an
impossibility.
\end{proof}

\section{Localizing a bad crossing}

\begin{lemma}\label{lem:supports}
Let $J_0\subset[N/2,4N]$ be a block of $B$ consecutive integers and put
$D_N=\lceil\ell^{11}\rceil$. On either side of each $x\in[N/2,4N]$,
within distance $D_N$, there is a composite integer $v$ such that
$\{v\}\times J_0\subset\cS$.
\end{lemma}
\begin{proof}
Let $Q_0$ be the set of all prime divisors of the integers in $J_0$.
Its cardinality $k$ satisfies
\[
                   k\leq B\frac{\log(4N)}{\log2}=O(\ell^3).
\]
Choose a prime $q\notin Q_0$ with $q\leq p_{k+1}$.
By~\eqref{eq:jacobsthal}, integers $a$ coprime to $\prod_{p\in Q_0}p$
occur with gaps $O((k+1)^2\log^4(k+2))$.
Consequently integers $v=qa$ of this kind occur with gaps
$O((k+1)^3\log^5(k+2))$.
For $v$ near the indicated $x$ we have $a>1$, because $q=N^{o(1)}$ and
$x\geq N/2$. Hence $v$ is composite. It is coprime to every element of
$J_0$, giving the required safe segment.
The gap bound is $O(\ell^9(\log\ell)^5)=o(\ell^{11})$; allowing an
additional endpoint gap yields a support strictly on either side within
$D_N$ for all sufficiently large $N$.
\end{proof}

\begin{lemma}\label{lem:localized}
A bad top-bottom star crossing of a rectangle
$I\times J\subset[N/2,4N]^2$ with $|J|\geq N^{9/10}$ contains points
\[
                    v_i=(x_0+a_i,y_0+b_i),\qquad 1\leq i\leq r,
\]
where $(x_0,y_0)$ is a point of that crossing,
$|a_i|,|b_i|\leq R_N:=2D_N+B$, and there are distinct primes $q_i>z$
dividing both coordinates of $v_i$.
\end{lemma}
\begin{proof}
Choose $J_0$ and $c_i$ by Lemma~\ref{lem:band}.
Orient the crossing from bottom to top. Take its first visit to the upper
row of $J_0$ and its last preceding visit to the lower row. The intervening
subpath stays inside the closed band: a visit above its upper row would
contradict the first choice, and a visit below its lower row would force a
later return to the lower row. Let $(x_0,y_0)$ be its first vertex.

The two safe support columns from Lemma~\ref{lem:supports} cannot be
crossed by this bad subpath. A star step changes the first coordinate by at
most one, and the subpath cannot pass around an endpoint while remaining
in the band. Its horizontal span is therefore at most $2D_N$.
It visits every selected row $c_i$; choose $v_i$ on that row.
Since $c_i$ is composite, badness at $v_i$ implies nonvisibility, and a
prime $q_i$ dividing the common divisor exceeds $z$.
The $q_i$ are distinct because the $c_i$ are pairwise coprime.
The displayed offset bounds follow from the horizontal span and band height.
\end{proof}

\section{The interpolation equation}

The height of an integer polynomial is the largest absolute value of its
coefficients. The following elementary interpolation step keeps track of
height as well as degree.

\begin{lemma}\label{lem:interpolation}
Let $d\geq1$, $r=\binom{d+2}{2}-1$, and $R\geq1$.
For any $r$ integer pairs $(a_i,b_i)$ with $|a_i|,|b_i|\leq R$ there is
a nonzero $H\in\ZZ[A,B]$ of total degree at most $d$ such that
\begin{equation}\label{eq:height}
 H(a_i,b_i)=0\quad(1\leq i\leq r),\qquad
                         \htp(H)\leq r!R^{dr}.
\end{equation}
If distinct primes $q_i$ divide both coordinates of
$(x_0+a_i,y_0+b_i)$, then
\begin{equation}\label{eq:evaluation-divisibility}
                         \prod_{i=1}^r q_i\mid H(-x_0,-y_0).
\end{equation}
\end{lemma}
\begin{proof}
Evaluate the $r+1$ monomials of total degree at most $d$ at the given
points. This gives an $r$ by $r+1$ integer matrix whose entries have
absolute value at most $R^d$.
Let its rank be $t\leq r$. Choose $t$ independent rows, $t$ columns
forming a nonsingular minor, and one further column. The signed minors of
the resulting $t$ by $t+1$ matrix give a nonzero integer kernel vector.
Extend it by zero on all other columns. It annihilates the chosen rows,
hence every row. Each coefficient is at most $t!R^{dt}\leq r!R^{dr}$.
This proves~\eqref{eq:height}, without a full-rank assumption.

Modulo $q_i$ the offset point equals $(-x_0,-y_0)$. Hence $q_i$ divides
$H(-x_0,-y_0)$. Distinctness gives~\eqref{eq:evaluation-divisibility}.
\end{proof}

\begin{proposition}\label{prop:algebraic-point}
Set $\mathcal H_N=\lceil\exp((\log\log N)^6)\rceil$.
Every bad top-bottom star crossing of a rectangle
$I\times J\subset[N/2,4N]^2$ with $|J|\geq N^{9/10}$ contains a point
$(x_0,y_0)$ satisfying $F(x_0,y_0)=0$ for some nonzero
$F\in\ZZ[X,Y]$ of total degree at most $d$ and height at most
$\mathcal H_N$. The same statement holds for a left-right bad star crossing
when $|I|\geq N^{9/10}$.
\end{proposition}
\begin{proof}
Apply Lemmas~\ref{lem:localized} and~\ref{lem:interpolation}.
Because $R_N=O(\ell^{11})$ and $d=O(L)$, $r=O(L^2)$, we have
\begin{equation}\label{eq:coefficient-size}
             \log(r!R_N^{dr})=O(L^4)<L^6
\end{equation}
eventually. Also
\[
 \log\prod_iq_i>\frac r u\ell,\qquad
 |H(-x_0,-y_0)|\leq(r+1)\,r!R_N^{dr}(4N)^d.
\]
The logarithm of the latter bound is $d\ell+O(L^4)$, whereas
$r/u=8u+O(1)$ and $d=4u+O(1)$.
Thus the prime product strictly exceeds the bound for a nonzero evaluation.
The divisibility in~\eqref{eq:evaluation-divisibility} forces
$H(-x_0,-y_0)=0$. Set $F(X,Y)=H(-X,-Y)$.

For a horizontal bad crossing, interchange the coordinate roles in the
band and support construction. The interpolation applies to the original
coordinate pair, or equivalently its resulting polynomial is transposed
back. The point still lies on the original crossing and the degree and
height bounds do not change.
\end{proof}

\section{A fixed direction avoiding the equations}

For dyadic $N=2^j$, define
\begin{equation}\label{eq:rho}
             \rho_N=\exp\left(-\frac{\log N}{20000\log\log N}\right).
\end{equation}
Consider all nonzero polynomials $f\in\ZZ[T]$ of degree at most $d$ and
height at most $\mathcal H_N$. Around the real part of each complex root
put an interval of radius $4\rho_N$, and let $E_N$ be their union.
Constant polynomials contribute no intervals. Writing $|E_N|$ for
Lebesgue measure, we have
\begin{equation}\label{eq:exception-measure}
 |E_N|\leq8\rho_Nd(2\mathcal H_N+1)^{d+1}.
\end{equation}
Its logarithm is bounded above by
\[
              -\frac{\log N}{20000\log\log N}+O((\log\log N)^7).
\]
Therefore $\sum_j|E_{2^j}|<\infty$. For example, its summands are
eventually bounded by $j^{-2}$. The measure of the set lying in infinitely
many $E_{2^j}$ is at most every tail of this convergent series, hence is zero.

Fix $\alpha\in(4/3,5/3)$ outside that null set. It lies outside $E_N$
for all sufficiently large dyadic $N$. Define the lattice corridor
\begin{equation}\label{eq:corridor}
 \mathcal U_N(\alpha)=
 \{(x,y)\in\ZZ^2\cap[N/2,4N]^2:|x/y-\alpha|\leq\rho_N\}.
\end{equation}

\begin{lemma}\label{lem:avoidance}
No point of $\mathcal U_N(\alpha)$ is a zero of a nonzero integer
polynomial of total degree at most $d$ and height at most
$\mathcal H_N$, for all sufficiently large dyadic $N$.
\end{lemma}
\begin{proof}
Suppose $F(x,y)=0$ at such a point. A nonzero constant cannot vanish, so
let $m\geq1$ be its total degree and put
$f(T)=\sum_{a+b=m}F_{ab}T^a$, where $F_{ab}$ is the coefficient of
$X^aY^b$. This is a nonzero integer polynomial of degree at most $d$
and height at most $\mathcal H_N$.
Since $1<x/y<2$ eventually and $y\geq N/2$, the lower-degree terms give
\[
 |f(x/y)|\leq\frac{2(r+1)2^d\mathcal H_N}{N}
                                                \leq N^{-1/2}
\]
for all sufficiently large $N$. The factor $2^d$ is harmless because
$d=O(\log\log N)$.
If $f$ is constant, this contradicts its being a nonzero integer.
Otherwise factor it over $\CC$. Its leading coefficient has absolute value
at least one, so one of its roots $\beta$ satisfies
\[
 |x/y-\beta|\leq N^{-1/(2\deg f)}
                   \leq N^{-1/(2d)}\leq\rho_N^2.
\]
The last inequality uses $d\leq5000\log\log N$ eventually.
It follows that $|\alpha-\Re\beta|\leq\rho_N+\rho_N^2<4\rho_N$,
contradicting $\alpha\notin E_N$.
\end{proof}

\begin{corollary}\label{cor:crossings}
Every integer rectangle $I\times J\subset\mathcal U_N(\alpha)$ has
a safe left-right crossing if $|J|\geq N^{9/10}$, and a safe top-bottom
crossing if $|I|\geq N^{9/10}$, for all sufficiently large dyadic $N$.
\end{corollary}
\begin{proof}
Failure of the first crossing gives a bad top-bottom star crossing by
Lemma~\ref{lem:planar}. Proposition~\ref{prop:algebraic-point} contradicts
Lemma~\ref{lem:avoidance}. Use the horizontal versions for the second claim.
\end{proof}

\section{Joining the scales}

\begin{proof}[Proof of Theorem~\ref{thm:directions}]
Fix a direction $\alpha$ as in the preceding section, and take all dyadic
scales sufficiently large for the preceding conclusions. Put
\[
 h_j=\left\lfloor\frac{2^j\rho_{2^j}}{100}\right\rfloor,\qquad
 s_j=\lfloor h_j/10\rfloor.
\]
We have $h_j/(2^j)^{9/10}\to\infty$.
For $t$ large, the derivative of $t/(20000\log t)$ lies between zero
and $1/2$. Equation~\eqref{eq:rho} therefore implies
\[
 \sqrt2<\frac{2N\rho_{2N}}{N\rho_N}<2
\]
eventually. In particular the integers $h_j$ are nondecreasing and
$h_{j+1}\leq3h_j$ for all sufficiently large $j$.

At scale $N=2^j$ choose the integer heights $N,N+s_j,N+2s_j,\ldots$
up to $2N$, adding the endpoint $2N$ if necessary. At each selected height
$t$ let $Q(j,t)$ be the integer square with centre
$(\lfloor\alpha t\rfloor,t)$ and half-side $h_j$.
Every point of this square satisfies
\[
 |x-\alpha y|\leq1+(1+\alpha)h_j,\qquad y\geq N-h_j.
\]
Thus the square lies in $\mathcal U_N(\alpha)$, with ratio error less
than $\rho_N/10$, eventually. Its coordinates lie in $[N/2,4N]^2$.
Choose a safe horizontal crossing of each square by
Corollary~\ref{cor:crossings}.

For two consecutive squares $Q,Q'$ at the same scale, their centres
differ by at most $h_j/10$ vertically and $h_j/4$ horizontally, including
the rounding of $\alpha t$. Form a bridge rectangle $T$ by intersecting
their horizontal intervals and taking the union of their vertical intervals.
Its horizontal side length is at least $7h_j/4$ and its vertical side
length is at most $21h_j/10$, apart from the added endpoint site.
The same ratio estimate, with $h_j$ replaced by $5h_j/4$, puts $T$
inside $\mathcal U_N(\alpha)$. Its horizontal side has more than
$N^{9/10}$ sites, so it has a safe vertical crossing.

This vertical crossing meets the chosen horizontal crossing of each of
$Q,Q'$. Indeed, restricting it between the lower and upper sides of $Q$
gives a vertical subcrossing in the rectangle with $Q$'s vertical interval
and $T$'s horizontal interval. Restricting the horizontal crossing of $Q$
to that horizontal interval gives the opposite crossing of the same
rectangle. Lemma~\ref{lem:planar} gives an intersection. The argument
for $Q'$ is identical. All horizontal crossings at this scale are connected.

At height $2N$, the last square at scale $j$ and the first at scale $j+1$
have the same centre and half-sides $h_j\leq h_{j+1}$.
Use a bridge whose horizontal interval is that of the smaller square and
whose vertical interval is that of the larger. It lies in the larger square,
hence in $\mathcal U_{2N}(\alpha)$. Since
$2h_j+1>(2N)^{9/10}$ eventually, it has a safe vertical crossing by
Corollary~\ref{cor:crossings}. The preceding intersection argument joins
the horizontal crossings at the two scales.

The union of all these chosen paths is one connected unbounded subgraph
of $\cS$. Fix a vertex on the first horizontal crossing. It reaches
arbitrarily large first coordinates in this same union. The tree of finite
simple paths from this fixed vertex is finitely branching and has nodes
of arbitrarily large depth; Koenig's lemma gives a ray in the union.
Each finite collection of the early squares and bridges has finitely many
vertices. On every sufficiently late one the coordinate ratio is within
$\rho_{2^j}$ of $\alpha$. Since a ray eventually leaves every finite set
and $\rho_{2^j}\to0$, its ratio tends to $\alpha$.
\end{proof}

\begin{corollary}\label{cor:two-components}
The graph on $\cS$ has an infinite component in each of the regions $x>y$
and $x<y$. Removing finitely many vertices leaves an infinite component
in each region.
\end{corollary}
\begin{proof}
The construction gives a ray in $x>y$; transpose its coordinates for the
other region. No nearest-neighbour path can change the sign of $x-y$
without visiting the excluded diagonal. A ray has only finitely many
vertices in a prescribed finite set, so an appropriate tail survives its
removal.
\end{proof}

\section*{AI-assisted research}
During the preparation of this work, the author used an author-developed
local implementation of the Proof Engine infrastructure (development
version 2.0) for AI-assisted research. The project-level research framework is described
in \emph{Proof Engine 2.0}~\cite{LiProofEngineTwo2026}.
See also \emph{Proof Engine Infrastructure} (1.0)~\cite{LiProofEngine2026}
for the complementary proof-development framework.
The implementation uses replaceable AI-agent and execution-harness components;
the present work used OpenAI Codex as its AI-agent component. The author
is responsible for the definitions, statements, proofs, citations, code,
and exposition.

\begin{thebibliography}{9}
\bibitem{Erdos1980}
P.~Erd\H{o}s, A survey of problems in combinatorial number theory,
\emph{Annals of Discrete Mathematics} \textbf{6} (1980), 89--115.
\href{https://doi.org/10.1016/S0167-5060(08)70697-6}{doi:10.1016/S0167-5060(08)70697-6}.
\bibitem{Vaughan1977}
R.~C.~Vaughan, On the order of magnitude of Jacobsthal's function,
\emph{Proceedings of the Edinburgh Mathematical Society} \textbf{20} (1977), 329--331.
\href{https://doi.org/10.1017/S0013091500026560}{doi:10.1017/S0013091500026560}.
\bibitem{LeFournLiuMartineau2025}
S.~Le Fourn, M.~Liu and S.~Martineau,
Percolative properties of the random coprime colouring,
\href{https://arxiv.org/abs/2509.08452}{arXiv:2509.08452} (2025).
\bibitem{RosserSchoenfeld1962}
J.~Barkley Rosser and L.~Schoenfeld,
Approximate formulas for some functions of prime numbers,
\emph{Illinois Journal of Mathematics} \textbf{6} (1962), 64--94.
\href{https://doi.org/10.1215/ijm/1255631807}{doi:10.1215/ijm/1255631807}.
\bibitem{Vardi1999}
I.~Vardi, Deterministic percolation,
\emph{Communications in Mathematical Physics} \textbf{207} (1999), 43--66.
\href{https://doi.org/10.1007/s002200050717}{doi:10.1007/s002200050717}.
\bibitem{LiProofEngineTwo2026}
A.~Chengyu Li, Proof Engine 2.0: An Evidence-Governed Method for Autonomous
Mathematical Research, Zenodo Working Paper (2026).
\href{https://doi.org/10.5281/zenodo.22346064}{doi:10.5281/zenodo.22346064}.
\bibitem{LiProofEngine2026}
A.~Chengyu Li, Proof Engine Infrastructure: A Fail-Closed Claim-Graph Method
for Accountable AI-Assisted Mathematical Research,
SSRN Scholarly Paper 7237460 (2026).
\href{https://doi.org/10.2139/ssrn.7237460}{doi:10.2139/ssrn.7237460}.
\end{thebibliography}
\medskip
\noindent\small First publicly available version: 6 September 2026,
\href{https://doi.org/10.5281/zenodo.22448689}{doi:10.5281/zenodo.22448689}.
\end{document}
